\documentclass[../DoAn.tex]{subfiles}
\begin{document}
\chapter*{Abstract}

The thesis "Design and Development of a Web Application for Recommending Weekend Entertainment Places for Children" focuses on supporting families in searching, comparing, and selecting suitable places for children. In practice, information about entertainment places may be distributed across different sources and may include location, suitable age range, ticket price, facilities, images, reviews, and ticket availability. This makes the decision-making process time-consuming for users.

The proposed approach is to build TheWeekend as a web-based system with a backend API, a MongoDB database, and an Android scanner application for ticket check-in. The system organizes internal data about places, categories, tags, reviews, favorites, bookings, payments, and electronic tickets. In addition to common search and detail-view functions, TheWeekend provides an internal-data-based recommendation mechanism and a chatbot that combines Gemini API with data retrieved from MongoDB.

The implemented functions include user registration, login, role-based access control, place search, place detail viewing, favorites, reviews with image submission, chatbot question answering, ticket booking, PayOS payment, QR ticket generation, administration features, and QR ticket check-in through the Android scanner application. The system also includes components for ticket management, payment tracking, and user notifications.

The system was manually tested in a local environment for the main web functions and the Android scanner workflow. However, the thesis does not include an official user survey, load testing, or performance testing. Therefore, it does not draw conclusions about user satisfaction, scalability, performance, or the quantitative quality of the recommendation mechanism and chatbot.
\end{document}
